\subsection{DET-16: the arithmetic cosine kernel and full-multiplicity reconstruction}
The scalar prime-sum endpoint supplies a precise positive-kernel route under
RH. Its cyclic Hilbert space alone loses repeated-zero multiplicities.
The atom masses retain those integers, however, and determine a
full-multiplicity positive realization of \(B\). This uses the standard
cyclic construction for a positive-definite kernel; the specific atomic
model and every normalization needed here are proved below. No novelty
claim for that construction is intended.

\paragraph{The arithmetic kernel and its exact positivity criterion.}
Use the actual \(G_\zeta\) and explicit correction \(b\) of
\fid{STEP6-02}, and define, without a zero list,
\[
 k(t)=e^{-|t|/2}\bigl(G_\zeta(e^{|t|})-b(e^{|t|})\bigr),
 \qquad t\in\R,\qquad K(u,v)=k(v-u).
\]
This is a continuous real even function, with
\[
 k(0)=2+\gamma_E-\log(4\pi)=2\lambda_1.
\]
The exact completed-zero expansion is
\begin{equation}\label{eq:arithkernelzero}
 k(t)=\sum_\rho
 \frac{N_\rho e^{(\rho-1/2)t}}{\rho(1-\rho)}
 =2\Tr\bigl(C\cos(tA)\bigr).
\end{equation}
Here the first sum is over distinct nontrivial zeros and \(N_\rho\)
is their multiplicity; the second uses the one-per-reflection-pair
operators of \fid{STEP6-01}. Before RH, the associated \(\gamma\)'s
may be complex. The series converges absolutely locally uniformly.

\begin{proposition}
The following conditions are equivalent:
\[
 \mathrm{RH};\qquad k\text{ is bounded on }\R;\qquad
 K\text{ is positive definite}.
\]
Positive definite means
\(\sum_{i,j}\overline{c_i}c_jk(t_j-t_i)\ge0\)
for every finite list of real \(t_i\) and complex \(c_i\).
Under these conditions there is a finite positive even measure
\begin{equation}\label{eq:arithkernelmeasure}
 \mu_K=\sum_{\gamma>0}
 \frac{m_\gamma}{\gamma^2+1/4}
       (\delta_\gamma+\delta_{-\gamma}),\qquad
 k(t)=\int_\R e^{ixt}\,d\mu_K(x),
\end{equation}
where this displayed identification uses the distinct positive
Xi-zero ordinates and their multiplicities. Its total mass is
\(k(0)\), and it has no atom at zero.
\end{proposition}
\begin{proof}
The arithmetic identity and all endpoint conventions are those proved in
\fid{STEP6-02} from the corrected explicit formula
\cite[corrected Lemma 2.1, pp.~1--2]{ramare}.
For \(t\ge0\), divide that formula by \(e^{t/2}\).
The functional equation pairs \(\rho\) with \(1-\rho\), so the
zero sum is even in \(t\); conjugation makes it real.
On a bounded \(t\)-interval each exponential is bounded by
\(e^{|t|/2}\), and
\(\sum_\rho N_\rho/|\rho(1-\rho)|<\infty\).
This proves \eqref{eq:arithkernelzero}, its convergence and continuity.

Under RH, \(\rho=1/2+i\gamma\) and
\(\rho(1-\rho)=\gamma^2+1/4\).
Pairing the positive and negative ordinates gives
\eqref{eq:arithkernelmeasure}. The same summability proves finite
mass and uniform absolute convergence on all real \(t\).
For a finite list,
\[
 \sum_{i,j}\overline{c_i}c_jk(t_j-t_i)
 =\int_\R\left|\sum_jc_je^{ixt_j}\right|^2d\mu_K(x)\ge0.
\]
Also \(|k(t)|\le\mu_K(\R)=k(0)\).
The classical \(\Xi(0)\ne0\) excludes the zero atom.

Conversely, positivity of the two-point Gram matrices at \(0,t\)
gives \(|k(t)|\le k(0)\), so positive type implies boundedness.
If merely \(|k(t)|\le M\), then, for every \(X\ge1\),
\[
 |G_\zeta(X)-b(X)|\le M\sqrt X.
\]
Since \(b\) is bounded, this supplies a global eventual lower bound
\(G_\zeta(X)\ge-K\sqrt X\) with finite \(K\).
The proved Landau lemma in \fid{STEP6-03}, with \(\theta=1/2\),
excludes zeros with real part greater than \(1/2\).
The functional equation excludes those with real part less than \(1/2\).
Thus boundedness implies RH. This proves all equivalences.
\end{proof}

\paragraph{The exact Weil test space.}
For \(t\in\R\), put
\[
 q_t(x)=e^{-t/2}\mathbf1_{(0,e^t]}(x),\qquad
 \mathcal Mq_t(s)=\frac{e^{(s-1/2)t}}s\quad(\Re s>0).
\]
On the finite complex span of these functions define the
conjugate-linear-first form
\[
 \mathcal W(f,h)=
 \sum_\rho N_\rho\,
 \overline{\mathcal Mf(1-\overline\rho)}\,\mathcal Mh(\rho).
\]
This sum is absolutely convergent, and directly
\begin{equation}\label{eq:arithkernelweilgram}
 \mathcal W(q_u,q_v)=k(v-u).
\end{equation}
Indeed the summand is
\(e^{(\rho-1/2)(v-u)}/[\rho(1-\rho)]\).
The reflection-conjugation symmetry of the zeros makes
\(\mathcal W\) Hermitian. On real functions its quadratic value is
the \(W_{\rm std}\) normalization of \fid{STEP6-01}; under RH it is
\(\sum_\rho N_\rho|\mathcal Mf(\rho)|^2\).
Thus positivity of this specific dilation-generated Weil form is
equivalent to RH by the preceding proposition.

For \(X=e^t\ge1\) and \(a=X^{\eta-1/2}\), the source test in
\fid{STEP6-01} is \(q_0\pm a q_{-t}\). Its quadratic value is exactly
\[
 k(0)(1+a^2)\pm2a\,k(t)
 =2\lambda_1(1+X^{2\eta-1})
 \pm2X^{\eta-1}(G_\zeta(X)-b(X)).
\]
This checks the sign, both factors of two, and the half-center.
It specifies a candidate Weil form; it does not identify the
historical undefined \(W\).

\paragraph{Recovering the atoms from the arithmetic function.}
Assume RH, or equivalently the positive-type condition just proved.
For every real \(\lambda\), define its coefficient by the actual limit
\begin{equation}\label{eq:arithkernelatom}
 a_K(\lambda)=
 \lim_{T\to\infty}\frac1{2T}
       \int_{-T}^{T}k(t)e^{-i\lambda t}\,dt.
\end{equation}
The limit exists and equals \(\mu_K(\{\lambda\})\). To prove this,
insert the finite measure in \eqref{eq:arithkernelmeasure} and use
Fubini on the finite time interval. The inner average is
\(\sin(T(x-\lambda))/(T(x-\lambda))\), with value one when
\(x=\lambda\). Its modulus is at most one, and it tends to zero
at every other \(x\). Dominated convergence proves the assertion.
The same argument applies to any finite positive measure with Fourier
transform \(k\). Such a measure must have these same atoms; their
masses already sum to \(k(0)\), its entire mass, so it has no
remaining part. Thus the positive measure is uniquely determined.

Write
\[
 \Gamma_K=\{\lambda>0:a_K(\lambda)>0\},\qquad
 m_K(\lambda)=(\lambda^2+1/4)a_K(\lambda)
 \quad(\lambda\in\Gamma_K).
\]
The exact explicit formula proves that \(m_K(\lambda)\) is a
positive integer and is the full zero multiplicity.
Integrality is a consequence of this arithmetic identity, not a
property of an arbitrary positive-definite kernel.
Both the frequencies and these integers have now been specified
from \(G_\zeta\) and \(b\), rather than by an input enumeration of zeros.
For \(\lambda>0\), the same coefficient has the entirely arithmetic form
\[
 a_K(\lambda)=
 \lim_{Y\to\infty}\frac1{\log Y}
 \int_1^Y
 \frac{G_\zeta(X)-b(X)}{X^{3/2}}
       \cos(\lambda\log X)\,dX.
\]
The change of variables is \(X=e^t\), using evenness of \(k\).
The recovered set is discrete, has no finite accumulation point,
and stays a positive distance from zero, by its identification with
the nonzero Xi ordinates.

\paragraph{The cyclic construction and its multiplicity limit.}
Complete the finite span of symbols \(e_t\), with
\(\langle e_u,e_v\rangle=k(v-u)\), after dividing by its null space.
Call the resulting Hilbert space \(\mathcal H_{\rm cyc}\).
Translations \(U_se_t=e_{t+s}\) are unitary, and \(e_0\) is cyclic.
This is the usual cyclic construction of a positive-definite function.
Here it has the explicit unitary model
\[
 \mathcal H_{\rm cyc}\simeq L^2(\mu_K),\qquad
 e_t(x)=e^{ixt},\quad U_sh(x)=e^{ixs}h(x),\quad e_0=1.
\]
\begin{proof}
The displayed map is isometric by the Fourier formula.
Its range is dense: for any atom \(\lambda\), the strong vector average
\[
 \frac1{2T}\int_{-T}^{T}e^{-i\lambda t}e_t\,dt
\]
lies in the closed span of the \(e_t\)'s and converges in
\(L^2(\mu_K)\) to \(\mathbf1_{\{\lambda\}}\).
This is the same sinc calculation, now with its square bounded
by one and dominated convergence for the finite measure.
Finite linear combinations of the atomic indicators are dense.
The vector integrals exist because \(t\mapsto e_t\) is norm
continuous, as follows by dominated convergence.
The same argument gives strong continuity of \(U_s\).
\end{proof}
Its selfadjoint generator, with \(U_s=e^{isA_{\rm cyc}}\), is
\[
 A_{\rm cyc}h(x)=xh(x),\qquad
 D(A_{\rm cyc})=\{h:\int x^2|h(x)|^2d\mu_K(x)<\infty\}.
\]
Selfadjointness follows directly by testing the adjoint on the
orthonormal atomic basis
\(\mathbf1_{\{\lambda\}}/\sqrt{a_K(\lambda)}\).
No differentiability of the cyclic vector is assumed.

There is exactly one dimension at each distinct signed atom,
regardless of its mass. On the positive-frequency subspace let
\(B_{\rm red}=A_{\rm cyc}^{-2}\). It is positive and trace class, but
\[
 \det(I-wB_{\rm red})
 =\prod_{\lambda\in\Gamma_K}(1-w/\lambda^2).
\]
This equals \(G(w)\) exactly when all the zeros are simple.
The reflection-even half of the full cyclic space gives the same
reduced determinant; the full signed cyclic space gives its square.
Thus the cyclic construction alone does not restore multiplicity.

In particular, atom mass is not an ordinary trace multiplicity.
On the full cyclic space, the diagonal operator
\(D_Kh(\lambda)=a_K(\lambda)h(\lambda)\) is trace class and
\[
 k(t)=\langle e_0,U_te_0\rangle
     =\Tr(D_KU_t).
\]
By contrast, the ordinary trace of a trace-class \(f(A_{\rm cyc})\)
is \(\sum_{\lambda}f(\lambda)\), with one term per distinct atom,
not the integral of \(f\) against \(\mu_K\).

\paragraph{The ordinary determinant with all multiplicities.}
The recovered integers give the explicit Hilbert space and operators
\[
 \mathcal H_K=\bigoplus_{\lambda\in\Gamma_K}\C^{m_K(\lambda)},
 \qquad A_K|_\lambda=\lambda I,\quad
 C_K|_\lambda=(\lambda^2+1/4)^{-1}I,\quad
 B_K|_\lambda=\lambda^{-2}I.
\]
Take the maximal diagonal domain for \(A_K\). The spectrum with its
fiber dimensions is determined by the arithmetic kernel; its Hilbert
realization is canonical up to unitary choices inside these fibers.
This is a multiplicity reconstruction from the masses, not the cyclic
GNS space itself.

\begin{proposition}
Under RH, \(B_K\) and \(C_K\) are positive trace-class operators and
\[
 \Tr C_K=\frac{k(0)}2=\lambda_1,\qquad
 k(t)=2\Tr(C_K\cos(tA_K)),\qquad
 B_K=C_K(I-C_K/4)^{-1}.
\]
Their ordinary determinants satisfy
\[
 \det(I-wB_K)=\frac{\Xi(\sqrt w)}{\Xi(0)}=G(w),
 \qquad
 \det(I-s(1-s)C_K)=2\xi(s).
\]
\end{proposition}
\begin{proof}
Let \(\lambda_0=\inf\Gamma_K>0\). The ordinary eigenvalue sums are
\[
 \sum_{\lambda>0}\frac{m_K(\lambda)}{\lambda^2+1/4}
 =\sum_{\lambda>0}a_K(\lambda)=\frac{k(0)}2,
\]
\[
 \sum_{\lambda>0}\frac{m_K(\lambda)}{\lambda^2}
 =\sum_{\lambda>0}a_K(\lambda)
       \left(1+\frac1{4\lambda^2}\right)
 \le\left(1+\frac1{4\lambda_0^2}\right)\frac{k(0)}2.
\]
For these positive block-diagonal operators these are exactly their
trace norms, proving trace class. Absolute trace convergence gives
the cosine identity. Also
\(\|C_K\|\le(\lambda_0^2+1/4)^{-1}<4\), so the indicated inverse
is bounded, and its scalar action is exactly \(\lambda^{-2}\).

The atom identification gives the full Xi multiplicities. The
genus-zero product of \fid{B-01} therefore yields the first
ordinary determinant. Alternatively this identifies \(B_K\)
unitarily with the designated \(B\), without choosing a zero list
in the definition of \(B_K\).
For \(v\in\C\), the scalar spectral transform gives
\[
 \det(I-vC_K)=\frac{G(v-1/4)}{G(-1/4)}.
\]
Setting \(v=s(1-s)\), evenness of \(\Xi\) and
\(\xi(0)=1/2\) give \(2\xi(s)\), exactly as in \fid{DET-11}.
Thus the normalization of its resolvent trace is unchanged.
\end{proof}

One may also retain both signs, putting a copy of
\(\C^{m_K(|\lambda|)}\) at every signed atom. The vector with
coordinates \((\lambda^2+1/4)^{-1/2}\) in each such fiber has
matrix coefficient \(k(t)\). Its cyclic subspace has only the
constant-vector line in each fiber. For the inverse-square operator,
taking one sign or the even half under exchange of the signed fibers
gives a unitary copy of \(B_K\) and the single Xi determinant;
retaining both signs gives \(G^2\). The signed generator itself
exchanges the even and odd subspaces.

\paragraph{The resulting weights on the actual Bessel model.}
Retain every source hypothesis and the fixed orthogonal convolution
projections \(P_j\) of \fid{DET-15}. Enumerate the pairs
\[
 (\lambda,r),\qquad \lambda\in\Gamma_K,\quad
 1\le r\le m_K(\lambda),
\]
in increasing \(\lambda\) and then in increasing \(r\), and assign
\(\beta_j=\lambda^{-2}\) to that pair. These summable positive
weights are now prescribed by the scalar endpoint through
\eqref{eq:arithkernelatom}. Equivalently, first split the positive
atom into \(m=m_K(\lambda)\) equal weights
\[
 c=\frac{a_K(\lambda)}m=\frac1{\lambda^2+1/4},
 \qquad
 \beta=\frac{c}{1-c/4}
       =\frac{a_K(\lambda)}{m-a_K(\lambda)/4}
       =\lambda^{-2}.
\]
Using a single atom mass as a single projection weight would
give a different determinant.

The source convolution sums
\[
 T_N=\sum_{j\le N}\beta_jP_j,\qquad
 T_K^{\rm Bessel}=\sum_{j\ge1}\beta_jP_j
\]
therefore satisfy
\[
 T_N\ge0,\qquad
 \|T_K^{\rm Bessel}-T_N\|_1=\sum_{j>N}\beta_j,\qquad
 \det(I-wT_K^{\rm Bessel})=G(w).
\]
Indeed \fid{DET-15} proves that each finite sum is the bounded
extension of the actual convolution with
\(\sum_{j\le N}\beta_je_{u_j}\). The isometry sending the standard
fiber basis of \(\mathcal H_K\) to the corresponding orthonormal
Bessel vectors makes the limit \(B_K\oplus0\).
The measurable kernel and finite determinant convergence are
those established there. The embedding depends on the chosen
source vectors; no canonical embedding into the Bessel model
is asserted.

This closes a scalar-endpoint reconstruction under the stipulated RH:
the arithmetic function defines the kernel, its coefficients recover
the weights and multiplicities, and the ordinary determinant follows.
Its positivity is itself RH-equivalent, so it is not an independent
proof of RH. The reconstruction uses infinite Fourier averages and
the classical explicit formula to identify their spectral meaning.
It does not prove that the historical finite period matrices,
their signed coefficients, or the undefined \(W\) are these operators.
The actual source's weighted-character and power-window quantifiers
remain distinct from the global bound used in the converse.
In the Grothendieck comparison, this is still an analytic
trace-class and exterior-power construction, not a new arithmetic
cohomology or Frobenius realization.
